\section{Introduction}
\label{sec:intro}

Mixture-of-Experts (MoE) improves nominal model capacity by activating only a small subset of parameters for each token, but the scientific foundation of MoE is not sparsity by itself. It is \emph{expert specialization}. An MoE becomes genuinely useful only when different experts learn meaningfully different transformations, respond to meaningfully different parts of the data distribution, and together convert architectural multiplicity into \emph{effective} capacity. If this does not happen, additional experts remain largely redundant modules: the model pays the systems cost of multiple experts without obtaining the intended functional diversity \citep{fedus2022switch}. From this viewpoint, expert specialization is not an auxiliary property of MoE; it is the condition under which MoE works at all.

In practice, however, expert specialization is often much weaker than the architecture seems to suggest. Experts can look balanced in usage yet remain functionally similar; they can appear dispersed in parameter space yet still absorb highly correlated learning signals; and they can process superficially different tokens while relying on the same underlying representational directions. \textcolor{red}{Here we should insert the finalized empirical evidence from our main-model audit and probe experiments, showing exactly where current experts fail to specialize and which commonly used proxies break down.} The important point is that current failure is not exhausted by occupancy imbalance, expert under-utilization, or a single overlap statistic. The deeper problem is that the system often fails to produce experts that are simultaneously different in parameters, different in the data manifolds they process, and different in the updates they accumulate.

This also clarifies the limitation of many existing lines of work. Standard MoE training specifies specialization mostly \emph{indirectly}: sparse activation decides which experts are active, load-balancing objectives regulate how often each expert is used, shared-expert or fallback designs absorb common knowledge, and various regularizers encourage experts to look different in parameter space. These mechanisms can improve utilization, stability, or apparent diversity, but they do not directly enforce the stronger condition we actually care about: that the gate and the experts cooperate to form a system in which different experts consistently absorb different semantic submanifolds and therefore evolve along different learning trajectories. Parameter independence by itself is not enough; occupancy balance by itself is not enough; even geometric separation by itself is not enough. A more essential notion of specialization must jointly involve \emph{which data manifold is routed where} and \emph{how that routing reshapes expert-specific updates}.

Our current analysis is anchored in the geometry of hidden representations, especially mean bias and anisotropy \citep{ethayarajh2019contextual,chi2022representation}. Language-model activations contain a strong shared/common component, and this common structure can dominate similarity relationships before any expert-specific signal becomes visible. As a result, tokens that should be separated for specialization may first appear close along shared, high-energy directions, while semantically discriminative structure is relegated to weaker residual components. This makes representational geometry a natural starting point: if the input to the expert system is already dominated by a common mode, then apparent diversity at the architectural level need not translate into true specialization at the learning level.

The key lesson of our current milestone is therefore negative but important: making representations look more separable does not automatically make experts more specialized. Our recent experiments suggest that the missing link is the coupling between representational separation and expert-conditioned gradient differentiation. \textcolor{red}{This is where the final experimental section should show, with finalized metrics and ablations, that cleaner routing geometry or better-looking separation does not automatically imply stronger expert specialization under a stricter criterion.} This reframes the main question of the paper. Rather than starting from a fixed routing design and asking whether it works, we ask a more foundational question: \emph{what should count as expert specialization, and what kind of gate--expert system can actually realize it?}

Accordingly, the goal of this work is to move from superficial proxies of specialization toward a more intrinsic measure. In the current draft, we treat expert specialization as a joint property of three coupled objects: expert parameters, the routed data submanifolds, and the expert-conditioned updates induced by those routed data. This viewpoint is directly motivated by our empirical pipeline and will guide the remainder of the paper. \textcolor{red}{The final method is intentionally left open at this stage. The present draft should first establish the problem formulation, the failure diagnosis, and the evaluation target; a definitive routing architecture should be introduced only after the corresponding evidence becomes mature.}

At the current checkpoint, the introduction is organized around the following claims:
\begin{itemize}
    \item \textbf{Expert specialization is a coupled systems property.} It cannot be reduced to parameter overlap, expert load, or any single occupancy statistic.
    \item \textbf{A stronger notion of specialization must connect experts to data manifolds and updates.} Experts are specialized only when the routed data they receive and the gradients they accumulate both become differentiated.
    \item \textbf{Mean bias and anisotropy provide the current diagnostic anchor.} The representational common mode is a plausible upstream source of weak specialization, but \textcolor{red}{the final operational metric and the final routing mechanism remain to be finalized by the next round of experiments.}
\end{itemize}

In this sense, the contribution of the current draft is to sharpen the problem itself: what MoE needs is not merely a more elaborate expert inventory, but a principled account of how specialization is created, measured, and sustained.
